\section{Discussion}

Although the three main contributions of this thesis address challenges at different stages of the life cycle of a docked \ac{bss}, they share common underlying mechanisms. This section discusses these shared insights and reflects on their relevance for research and operational practice in bike-sharing and data-driven urban mobility.

A central theme common to all three contributions is the role of high-resolution and high-quality data as essential for meaningful analysis and prediction in \acp{bss}. Across system planning and operations, the results show that detailed spatiotemporal, contextual, and operational data enable the identification of clear heterogeneities within the system. Different parts of the city, stations, and bicycles exhibit distinct patterns of accessibility, usage, energy depletion, and component wear, which would be largely obscured under aggregated representations. Moreover, beyond enhancing descriptive analyses, high-resolution data can substantially improve modeling and prediction by reducing the need for simplifying assumptions.

Another important topic concerns the role of network effects in shaping system behavior. The analyses consistently indicate that outcomes depend not only on the local attributes of stations or trips, but also on their position within the broader cycling network and their relationships to other system elements. Inter-station proximity and system accessibility influence feasible station layouts; mobility flows are constrained by the structure of the cycling network rather than by simple distance alone; and repeated traversal of specific network segments contributes to uneven battery usage and component degradation. These observations underscore the need to treat \acp{bss} as networked systems, where spatial structure and connectivity determine how local characteristics translate into system-wide patterns. To capture these effects realistically, modeling approaches should go beyond distance-based representations and also account for topography and direction-dependent travel costs.

Within this networked context, structural design choices—particularly those related to station placement—play a key role in shaping operational dynamics. While this thesis does not explicitly analyze the relationship between planning decisions and operational outcomes, the results suggest that station layout influences mobility patterns, which in turn affect battery dynamics and maintenance needs. This perspective highlights a limitation of approaches that address planning, rebalancing, or maintenance in isolation and points to the potential value of integrated frameworks that consider these components as interconnected parts of a single system. At the same time, user behavior remains inherently adaptive, suggesting that data-driven models should be continuously updated to reflect evolving usage patterns and responses to system changes.

In parallel, ongoing technological developments are expanding both the scope and the practical relevance of data-driven approaches in \acp{bss}. New generations of shared bicycles increasingly incorporate \ac{gps} units and a growing range of onboard sensors, enabling continuous monitoring of location, usage, and mechanical conditions, including signals such as battery status, tire pressure, and vibration that provide additional information on the state of the bicycle. At the same time, the availability of contextual information is steadily improving, driven by the expansion of open urban data and more widespread data-sharing practices among public authorities. In this context, operators are becoming increasingly open to the adoption of \ac{ai}-based methods as a means to improve efficiency, reduce operational costs, and optimize resource allocation. As a result, although some of the analyses developed in this thesis—particularly those related to battery dynamics and predictive maintenance—remain relatively uncommon in the current literature, the increasing availability of rich operational and contextual data is likely to make the application of such methods progressively more feasible in real-world settings.

Beyond methodological insights, the findings of this thesis have implications for both \ac{bss} users and operators, as well as for society more broadly. While the analyses presented in this thesis are grounded in Barcelona’s \ac{bss}, the methodological approaches are largely transferable to other urban contexts, provided that comparable spatial, operational, and contextual data are available. For users, improved station placement and more reliable vehicle availability can translate into higher and more consistent service quality. In particular, planning approaches that explicitly account for accessibility and equity considerations can support a more balanced spatial coverage, helping to extend the benefits of bike-sharing to a wider range of neighborhoods and user groups. For operators, data-driven planning and predictive tools offer opportunities to allocate resources more efficiently and to reduce operational costs. At a societal level, more efficient, equitable, and reliable \acp{bss} can strengthen their role as a sustainable urban mobility option, supporting modal shift away from private motorized transport and contributing to lower emissions, reduced congestion, and more livable cities. In this sense, the approaches presented in this thesis align operational efficiency with broader environmental and social goals.

